\subsection{Marked traces and the actual generator relations}
The trace-only construction leaves open the arithmetic
comparison for specified operators. A regular central
insertion already detects the collision in pure or
integrally mixed coefficient models. Requiring composition
for just two generator words gives an obstruction without
any purity or algebraicity assumption.

\begin{finding}[\fid{F54}, Part 3: marked-trace comparison gates]
\label{f:markedtracegates}
Let $b>1$, and use the actual rank-$64$ tensor module $M$
and named generators $P_1,P_2$.
Put
\[
 Z=\frac{P_1+P_2}{D}
       -\frac{P_1P_2+P_2P_1}{D^2},\qquad
 z_1=\frac{D-1}{D},\quad z_2=\frac{4(D-1)}{D^2},
\]
\begin{equation}\label{eq:actualmarkedX}
 X=\frac{Z(z_1-Z)}{z_2}.
\end{equation}
This is a regular central operator at $v=1$.
It acts by $\Delta$ on $N(D-2)^4$, a rank-$8$ sector,
and by zero on the other sectors; in particular $X(1)=0$.

Suppose $F$ is an invertible algebraic matrix family on $M$,
regular at a place above $v=1$, commuting with the actual
action, and with regular evaluations at $b$ and all
sufficiently large $b^n$.
The marked trace identity
\begin{equation}\label{eq:markedXretuning}
 \Tr(X_{b^n}F_{b^n})=\Tr(X_bF_b^n)
\end{equation}
is impossible if the eight base eigenvalues on the
$\Delta$-sector either have one common absolute value
or have several integral weights
$|\alpha_j|=b^{w_j}$, $w_j\in\Z$.
No full-centralizer comparison is needed for this assertion.

Alternatively, suppose $F$ acts on a regular auxiliary
family $H$ whose full relative centralizer is the actual
order, with the insertion $X$ corresponding to its
specified central coordinate, and with the same algebraicity,
invertibility and regular-evaluation hypotheses.
Then \eqref{eq:markedXretuning} is impossible without
any weight restriction.

Finally, let $M_n$ be the actual module at $v=b^n$, let
$P_n$ be its named $P_1$, and let $F_n$ be arbitrary
invertible operators, with $F_1P_1=P_1F_1$ at the base.
The two identities
\begin{equation}\label{eq:twowordretuning}
 \Tr(P_n^jF_n)=\Tr(P_1^jF_1^n),\qquad j=1,2,
\end{equation}
cannot both hold for every sufficiently large $n$.
Here the right insertions are fixed actual base operators,
and the second is genuinely the square of the first.
This last assertion needs no algebraicity, purity,
full-centralizer recovery, or collision specialization.
\end{finding}

\begin{proof}
\par\smallskip\noindent\textbf{The marked central element.}\quad
In the actual algebra set $e=P_1/D$, $f=P_2/D$.
These are idempotents, and
\[
 Z=e+f-ef-fe.
\]
A direct multiplication gives
\[
 Ze=eZ=e-efe,\qquad Zf=fZ=f-fef.
\]
Thus $Z$ and $X$ are central.
In the row-action convention, the actual matrices on $N(a)$
are
\[
 P_1=\begin{pmatrix}D&0\\a&0\end{pmatrix},\qquad
 P_2=\begin{pmatrix}0&a\\0&D\end{pmatrix}.
\]
They give
\[
 Z=(1-a^2/D^2)I.
\]
On both scalar sectors $Z=0$.
Since $Z=z_1$ on $N(\delta)$ and $Z=z_2$ on $N(D-2)$,
\begin{equation}\label{eq:actualXbranches}
 X|_{\epsilon_\pm}=0,\qquad
 X|_{N(\delta)}=0,\qquad
 X|_{N(D-2)}=\Delta I.
\end{equation}
The denominators $D,z_2$ are units at $v=1$ and on
the previously specified base.
Equivalently, if
\[
 B=D(P_1+P_2)-(P_1P_2+P_2P_1),
\]
then
\[
 X=\frac{B((D-1)D-B)}{4(D-1)D^2}.
\]
This displays a regular rational expression in the named
generators. It is not an unspecified observable.

Under the actual coherent algebra model, $X$ is zero
on both scalar factors and the coordinate $x$ on
$M_2(S[x]/(x(x-\Delta)))$.
On the actual module, the regular \fid{F21} splitting
makes \eqref{eq:actualXbranches} an integral action identity.
Its $\Delta$-sector has rank $2\cdot4=8$, and $X(1)=0$
as a matrix, not merely at the level of eigenvalues.

\par\smallskip\noindent\textbf{Three scalar consequences.}\quad
Let $g$ be one algebraic function over $\C(v)$,
evaluated at $b^n$ through places of that same relation.

First, if
\[
 g_{b^n}=C\alpha^n
\]
eventually, where $C,\alpha\ne0$, then for some
$L>0$, $m\in\Z$,
\begin{equation}\label{eq:singleexponentialalgebraic}
 g(v)^L=C^L v^m.
\end{equation}
Indeed, expand its irreducible polynomial relation
$P(b^n,C\alpha^n)=0$.
This is a finite power sum on roots $b^i\alpha^j$.
Vandermonde uniqueness forces two distinct monomials
to have the same root. Their $j$ exponents differ,
so $\alpha^L=b^m$ for some $L,m$.
On $n=Lr$, substitute $v=t^L$ to obtain
$P(t^L,Ct^m)=0$ at infinitely many $t=b^r$.
It is therefore a rational identity.
The irreducible polynomial of $g$ has the root $Ct^m$,
so it divides $Y^L-C^Lv^m$ over $\C(v)$.
This proves \eqref{eq:singleexponentialalgebraic}.
In particular $g$ cannot vanish at any place over $v=1$.

Second, if
\[
 g_{b^n}=\sum_{j=1}^r\alpha_j^n
\]
with genuine multiplicities and a common modulus,
Part~1 gives a common $L>0$ and one integer $m$ such that
$\alpha_j^L=b^m$ for all $j$.
At degrees $Lr'$ the trace is $r b^{mr'}$.
Hence the algebraic function $g^L-r^L v^m$ vanishes
at places over infinitely many distinct parameters.
The nonzero-constant-term argument for its minimal
polynomial, proved in Part~1, forces
\begin{equation}\label{eq:markedcommonradius}
 g(v)^L=r^L v^m.
\end{equation}
Again $g$ cannot vanish above $1$.
This consequence uses no integral-weight assumption
and no continuity of the evaluation sheets.

Third, suppose the base roots have integral weights:
\[
 \alpha_j=b^{w_j}\zeta_j,\qquad
 w_j\in\Z,\quad |\zeta_j|=1.
\]
The mixed-integral trace argument of \fid{F51} gives
\begin{equation}\label{eq:markedintegraltrace}
 g(v)=\sum_{j=1}^r v^{w_j},\qquad g(1)=r.
\end{equation}
To recall the exact deduction, expand
$P(b^n,\sum_j b^{w_jn}\zeta_j^n)$ and group by the
integer power of $b^n$.
Each remaining coefficient is a finite power sum
on unit-modulus roots.
Distinct radii and Vandermonde uniqueness make every
coefficient sequence zero, including at $n=0$.
Those identities are the Laurent coefficients of
$P(v,\sum_jv^{w_j})=0$.
Irreducibility over $\C(v)$ makes this rational root
the algebraic function itself.
Negative weights are allowed in $\C(v)$.
These arguments keep genuine positive multiplicities;
they do not restore eigenvalues canceled in a virtual trace.

\par\smallskip\noindent\textbf{The actual tensor-module gate.}\quad
Since $F_b$ commutes with the actual action,
the rank-$8$ image of $X_b/\Delta(b)$ is invariant
and its restriction is invertible.
Let its eigenvalues be $\alpha_1,\ldots,\alpha_8$,
including algebraic multiplicities.
The algebraic function
\[
 g(v)=\Delta(b)^{-1}\Tr_M(X(v)F(v))
\]
would satisfy
\[
 g_{b^n}=\sum_{j=1}^8\alpha_j^n
\]
under \eqref{eq:markedXretuning}.
But regularity and $X(1)=0$ give $g(1)=0$.
Equation \eqref{eq:markedcommonradius} contradicts this
if the branch spectrum has one common modulus;
\eqref{eq:markedintegraltrace} gives the contradiction
$g(1)=8$ in the integral-mixture case.

A pure Weil spectrum supplies the common-radius hypothesis
at any chosen complex embedding.
No full point-Weil endomorphism algebra has been identified
with the actual algebra in this argument.
The same reasoning works when $X(1)$ is only nilpotent,
provided it commutes with $F(1)$:
their product is then nilpotent and its trace is zero.
For unrestricted multiple radii with nonintegral weights,
neither scalar conclusion is asserted.
For example, on $v=t^2$ the function $g=1+t$ has values
$1+(\sqrt b)^n$ on the positive sheets over $b^n$,
but vanishes at $t=-1$ over $v=1$.
This is a scalar boundary example, not an actual-action
or full-order realization.

\par\smallskip\noindent\textbf{The full relative-centralizer gate.}\quad
Now suppose $\End_F(H)$ is the actual regular order
and the central marking is preserved.
The operator $F$ is central in its full centralizer,
so its component in the matrix-center factor is
$c_0+c_1x$.
Its value on the $x=\Delta$ branch is
$d=c_0+c_1\Delta$.

Over the generic field the actual algebra is split
semisimple, with two scalar and two two-by-two matrix
factors.
A central element is a tuple of scalars, so $F$ is
diagonalizable.
Distinct primitive central blocks have distinct eigenvalues;
otherwise maps between equal-eigenvalue blocks would
enlarge the full centralizer.
Each two-by-two factor is thus the full endomorphism
algebra of a two-dimensional eigenspace.
It follows that
\[
 g(v)=\Tr_H(X(v)F(v))=2\Delta(v)d(v).
\]
At the base,
\[
 \Tr_H(X_bF_b^n)=2\Delta(b)d(b)^n.
\]
The constant $2\Delta(b)$ and the eigenvalue $d(b)$
are nonzero.
If \eqref{eq:markedXretuning} held, the single-exponential
identity \eqref{eq:singleexponentialalgebraic} would make
$g$ a unit at every point over $1$.
Its regular expression $2\Delta d$ instead vanishes there.
This proves the assertion without any weight assumption.

The trace-only model of Part~2 has
$g(v)=2\Delta(v)d(v)$ and $d(2)=8$.
The extra comparison would require
$g(2^n)=2\Delta(2)2^{3n}$.
It fails already at $v=4$, where the exact difference is
\[
 -\frac{197896315512672}{4359848400625}.
\]
Thus that model gives a concrete failure of the marked
axiom while preserving both identity traces.

\par\smallskip\noindent\textbf{Two actual inserted words.}\quad
Finally use the actual relation
\[
 P_n^2=D(b^n)P_n
\]
and the fixed base relation $P_1^2=D(b)P_1$.
If both identities in \eqref{eq:twowordretuning} held,
they would imply
\[
 \bigl(D(b^n)-D(b)\bigr)\Tr(P_1F_1^n)=0.
\]
The function $(v+v^{-1})^2$ is strictly increasing
for $v>1$, so the first factor is nonzero for $n>1$.
Hence $\Tr(P_1F_1^n)$ must vanish eventually.

But $P_1/D(b)$ is an idempotent of rank $40$:
the sectors contribute $20+16+4$.
Its image is $F_1$-invariant because the two commute,
and the restricted $F_1$ is invertible.
If its eigenvalues are $\lambda_i$, with genuine
multiplicities, then
\[
 \Tr(P_1F_1^n)=D(b)\sum_i\lambda_i^n.
\]
Grouped distinct nonzero roots have positive integer
coefficients.
Vandermonde uniqueness forbids their vanishing for
every sufficiently large $n$.
This is the required contradiction.
No commutation hypothesis at the other parameters
was needed for this last calculation.
\end{proof}

\paragraph{The comparison axiom is part of the conclusion.}
The left insertion in \eqref{eq:markedXretuning} is the
actual $X$ at $b^n$; the right insertion is the fixed
$X_b$.
Its trace law is additional data and does not follow
from identity-trace equality.
Although $X$ is a regular expression in named generators,
its parameter-dependent rational coefficients cannot
silently be moved through a fixed-operator comparison.

The two-word obstruction strengthens \fid{F42} to a
finite test under an explicit action and composition
requirement.
It does not exclude a dictionary in which the formal
coefficient $D(v)$ is represented by a graded coefficient
object, and the quadratic relation becomes a tensor
relation with that object.
Such a dictionary does not impose the fixed scalar
$D(b)$ used on the right of \eqref{eq:twowordretuning}.
The separate word objects of \fid{F45}, the coefficient
kernels of \fid{F43}, and the fixed-parameter constructions
of \fid{F50} and \fid{F53} remain outside that obstruction.
The remaining comparison concerns the designated
arithmetic, grading and basis structures.

% Separate paragraph for Part III.
\paragraph{\fid{F54}, Part 3, in Grothendieck's language.}
The marked insertion belongs to the actual coherent
algebra action, rather than being reconstructed from an
unmarked scalar trace.
Its trace-retuning law is a specified additional
arithmetic comparison.
The operator/representation language is SGA~3, I,
4.7.2--3, with cote~90, scan~17 as the audited framework
anchor; the underlying coherent module language is
EGA~I, Theorem~1.3.7, p.~86, Corollary~1.3.8, p.~87,
and Theorem~1.4.1, pp.~90--92
\cite{SGA3,EGA,Archive}.
A changing coefficient object belongs to the tensor
and Grothendieck-group side of the comparison, as in
SGA~6, V, \S\S2--3 and cote~133, scan~18.
No eigenvalue or purity theorem is inferred merely
from the existence of a coherent module.
The finite-power, scalar-function and two-word
obstructions have all been proved explicitly above.

